\scanpage{2}
\mspage{1}
\margin{\unsure{Ch|Cf} I}

Vers une géométrie des formes (topologiques)

(\uline{Juin \unsure{86}})

\margin{(5 juin)}
Approche vers une construction récurrente (sur l'entier naturel $n$) d'une ``géométrie des formes de dimension $\leq n$''.

Une \deleted{\unsure{``espace''}} ``forme de dim 0'' sera par définition un ensemble discret dont les éléments sont appelés les ``lieux'' de la forme.

\margin{Un modèle d'une forme de dim 0 sera un ens. $L$ (``lieux'') avec une relation d'équivalence (``non disjoints''). La forme \unsure{du} modèle est l'ens quotient.}

\uline{\deleted{\unsure{Formes}}\inserted{Modèle} de dimension 1}

Un \deleted{tel} tel modèle comporte

\deleted{
\textcircled{1} Un ensemble des ``\uline{lieux}'' de la forme, $L$.

\textcircled{2} Un sous-ensemble $L_0$ de $L$, formé des lieux dits ``\uline{critiques}'' ou ``\uline{de branchement}''.

3) Une relation \inserted{symétrique} dans $L$ (i.e. $R\subset L\times L$) ``$x$ et $y$ sont \uline{disjoints}'', impliquant la relation ``$x$ et $y$ sont distincts''.

4) Une relation \inserted{symétrique} plus fine ``$x$ et $y$ sont \uline{adjacents}''

\textcircled{5} Une \deleted{ensemble de parties de $L$, (appelés \inserted{\unsure{comme support}} ``\uline{segments}'' \unsure{``$x$ est s}} \inserted{\unsure{S} relation plus fine encore} (\unsure{comme ci dessus} pas symétrique) ``$x$ \deleted{\unsure{est}} est spécialisation de $y$'' (qui implique que $x$ est critique, et $y$ n'est pas critique).

\textcircled{6} Un ensemble $S$, \deleted{\unsure{partition de $L$}, \inserted{dont les éléments, \unsure{``segments''}}} \unsure{appelés} ``\uuline{segments}''

\textcircled{7} Pour tout segment \deleted{\unsure{un}} $I$, deux parties $\unsure{\delta I|\partial I}$ de
}
\rotatebox{90}{\reflectbox{\unsure{EXTRA STRONG}}}

\rotatebox{90}{\reflectbox{\unsure{EXTRA STRONG}}}
\mspage{1}
